\documentclass[../main]{subfiles}
\setcounter{chapter}{9}
\begin{document}
\chapter{Deslizamiento y sincronismo}
\begin{center}
\img{images/08/logo.jpg}{Motor Trif\'asico}
\end{center}

\subsection{Sincronismo}
La velocidad de giro del campo magn\'etico{} se conoce como velocidad de sincronismo $n_s$. El motor s\'incrono{} gira a la misma velocidad que el campo magn\'etico.
\begin{equation}
n_s = f \frac{60}{p}
\end{equation}

donde:
\begin{itemize}
\item $n_s$: Velocidad de sincronismo en revoluciones por minuto \unit{rpm}.
\item $f$: Frecuencia de la red en Hertz \unit{Hz}.
\item $p$: N\'umero{} de pares de polos del motor.
\end{itemize}

\begin{table}[h!]
\centering
\begin{tabular}{|c|c|c|}
\hline
$N^o$ de Polos & 50 \unit{Hz} & 60 \unit{Hz} \\ \hline \hline
2 polos & 3000 & 3600 \\
4 polos & 1500 & 1800 \\
6 polos & 1000 & 1200 \\
8 polos & 750 & 900 \\
12 polos & 500 & 600 \\
16 polos & 375 & 450 \\
24 polos & 250 & 300 \\ \hline
\end{tabular}
\caption{Velocidad Sincr\'onica{} para motores a 50 y 60 Hz.}
\label{tab:vel_sincrona}
\end{table}

\subsection{Deslizamiento}

\web{https://www.youtube.com/watch?v=duRA-drcfGQ}{Principio de funcionamiento de motor asincr\'onico}
\info{Definici\'on{} de deslizamiento}{El \emph{deslizamiento o resbalamiento} en una m\'aquina{} el\'ectrica{} es la diferencia relativa entre la velocidad del campo magn\'etico{} (velocidad de sincronismo) y la velocidad del rotor. En motores asincr\'onicos{} esta velocidad no es igual.}

Las siguientes expresiones son equivalentes para hallar el deslizamiento:
\begin{equation}
s = \frac{\omega_s - \omega_m}{\omega_s} \cdot 100\% = \frac{n_s - n_m}{n_s} \cdot 100\%
\end{equation}

Donde:
\begin{itemize}
\item $s$: Velocidad de deslizamiento (expresada en porcentaje).
\item $\omega_{s}$: Velocidad angular de sincronismo en radianes por segundo \unit{rad/s}.
\item $\omega_{m}$: Velocidad angular del rotor en radianes por segundo \unit{rad/s}.
\item $n_{s}$: Velocidad angular de sincronismo en revoluciones por minuto \unit{rpm}.
\item $n_{m}$: Velocidad angular del rotor en revoluciones por minuto \unit{rpm}.
\end{itemize}

El deslizamiento es especialmente \'util{} cuando analizamos el funcionamiento del Motor asincr\'onico{} ya que estas velocidades son distintas. El voltaje inducido en el bobinado rot\'orico{} de un motor de inducci\'on{} depende de la velocidad relativa del rotor con relaci\'on{} a los campos magn\'eticos.
Es posible expresar la velocidad mec\'anica{} del eje del rotor, en t\'erminos{} de la velocidad de sincronismo (velocidad del campo magn\'etico) y del deslizamiento.
\begin{equation}
n_{m} = n_{s}(1-s)
\end{equation}

\info{Rango de deslizamiento}{
En el Motor asincr\'onico{} $0 < s < 1$ ($-1 < s < 0$ funcionando como generador, depende la aplicaci\'on{}, alrededor de \qty{5}{\%}). En el Motor s\'incrono{} $s = 0$.
}

\subsection{Par Motor y velocidad}
\web{https://www.youtube.com/watch?v=vcr5pkT6p1o}{Arranque motores de inducci\'on}
\img{images/08/curva}{Curva Par motor-velocidad}

La velocidad nominal siempre ser\'a{} inferior a la velocidad de sincronismo o velocidad del campo magn\'etico{} del estator, debido al deslizamiento.
El par resistente aumenta a medida que aumenta la velocidad y el par motor disminuye al aumentar la velocidad.
Alrededor de la velocidad nominal existe una zona estable en la que el motor puede oscilar en funci\'on{} del par resistente, variando ligeramente la velocidad. Si lo sacamos de esa zona el motor saldr\'a{} de sincronismo y se parar\'a.
La Intensidad de arranque es muy elevada, del orden de 8 veces superior a la Intensidad nominal.
De ah\'i{} la necesidad de utilizar sistemas de arranque que minimicen la Intensidad de arranque.
\begin{equation}
P = \omega \cdot M
\end{equation}

Donde:
\begin{itemize}
\item $P$: Potencia en Watts \unit{W}.
\item $\omega$: velocidad del motor en radianes por segundo \unit{rad/s}.
\item $M$: Par motor en Newton-metro \unit{N.m}.
\end{itemize}

Multiplicaremos por $\frac{2\pi}{60}$ para pasar de R.P.M. a radianes/segundo.

\actividad{Resolver}
\begin{enumerate}
\item Un motor as\'incrono{} trif\'asico{} de rotor en cortocircuito posee una velocidad de campo $n_s$ de \qty{3000}{rpm}. \textquestiondown{}Cu\'al{} ser\'a{} el deslizamiento del rotor a plena carga si se mide con un tac\'ometro{} una velocidad de \qty{2750}{rpm}?
\item Un motor asincr\'onico{} trif\'asico{} de 1 par de polos funciona con \qty{60}{Hz}. Si su velocidad nominal es de \qty{3500}{rpm} \textquestiondown{}Cu\'al{} es su deslizamiento?
\item El deslizamiento de un motor monof\'asico{} de 2 pares de polos es de \qty{8}{\%}. Calcular la velocidad de sincronismo y la velocidad del motor.
\item El deslizamiento de un motor es de \qty{10}{\%} y su velocidad \qty{1600}{rpm}. Si el motor tiene 2 pares de polos \textquestiondown{}Qu\'e{} frecuencia lo est\'a{} alimentando?
\item Calcular la velocidad de sincronismo de un motor que tiene una velocidad de \qty{2801}{rpm} y un deslizamiento del \qty{6}{\%}.
\item Calcular el deslizamiento de un motor de 8 polos que gira a \qty{800}{rpm} y funciona a \qty{60}{Hz}.
\item \textquestiondown{}Con qu\'e{} frecuencia ser\'a{} necesario alimentar un motor asincr\'onico{} trif\'asico{} de \qty{1}{HP} para que tenga una velocidad de \qty{1000}{rpm} si posee 2 pares de polos?
\item El motor as\'incrono{} de \qty{1.5}{HP} gira a \qty{1400}{rpm}. \textquestiondown{}Cu\'al{} es el deslizamiento si es de \qty{50}{Hz} y cu\'antos{} pares de polos tiene?
\item El motor trif\'asico{} \emph{Marelli} de \qty{4}{HP} gira a \qty{2890}{rpm}. \textquestiondown{}Cu\'al{} es el deslizamiento si es de \qty{50}{Hz} y cu\'antos{} pares de polos tiene?
\item El motor monof\'asico{} asincr\'onico{} \emph{Industria Argentina} de \qty{1450}{rpm} de \qty{1}{HP}. \textquestiondown{}Qu\'e{} deslizamiento posee y cu\'antos{} pares de polos tiene?
\item Si los motores asincr\'onicos{} monof\'asicos{} normalizados tienen \qty{1400}{rpm} \textquestiondown{}Cu\'al{} es el deslizamiento para un motor normalizado?
\item En la placa de un motor asincr\'onico{} monof\'asico{} de \emph{Industria Argentina} figuran los siguientes datos:
\begin{itemize}
\item \qty{1}{HP}
\item \qty{1425}{rpm}
\item \qty{50}{Hz}
\item \qty{5.7}{A}
\item \qty{220}{V}
\item Cap \qtyrange{430}{514}{\micro F}
\item $\cos(\varphi) = 0,9$
\item IP21 uso general.
\end{itemize}
\textquestiondown{}Cu\'antos{} pares de polos tiene y cu\'al{} es su rendimiento?
\item En una publicaci\'on{} se vende un motor trif\'asico{} de velocidad \qty{3000}{rpm}, marca \emph{WEG}. Si posee un \qty{5}{\%} de deslizamiento \textquestiondown{}Cu\'al{} es la velocidad del motor a carga nominal?
\item Se determin\'o{} que un motor da 24 giros en 1 segundo. Adem\'as{}, el motor se alimenta con \qty{220}{V} y \qty{50}{Hz}. \textquestiondown{}Cu\'antos{} polos tiene el motor y cu\'al{} es su deslizamiento?
\item El eje de un motor tarda aproximadamente \qty{20,8}{ms} en dar una vuelta. Determinar si se alimenta con \qty{50}{Hz}, el n\'umero{} de pares de polos, la velocidad de sincronismo y la velocidad del motor en rpm.
\item Si un motor tiene un deslizamiento del \qty{8}{\%} y 2 pares de polos, alimentado por \qty{220}{V} y \qty{50}{Hz}. Determinar en cu\'antos{} milisegundos da la vuelta el eje.
\end{enumerate}
\end{document}